\chapter{Абсолютные условия \texorpdfstring{$A_1$--$A_{16}$}{A1-A16}}
\label{ch:axioms}

\section{Реестр условий}

Условия $A_1$--$A_{16}$ ограничивают допустимый локальный закон $g$ (\cref{ch:evolution}).
Ниже --- формулировки; следствия для $g$ выводятся в \cref{ch:evolution,ch:matter}.

\begin{axiom}[A1 · Каузальность]
Один тик $\le \hl$. $\varepsilon$ --- равные light-like NN; Moore и вторая оболочка запрещены.
$(3{+}1)$: $N_{12}$ FCC (\cref{ch:carrier}). $\ell_P:=\hl$.
\end{axiom}

\begin{axiom}[A2 · Локальность]
$g(x)$ зависит только от $\varepsilon$-окрестности сайта $x$.
\end{axiom}

\begin{axiom}[A3 · Унитарность]
$\sum_x |z(x,t)|^2$ invariant. Линейная часть --- product локальных unitaries;
нелинейность --- только $e^{i\Phi}$ на спиноре.
\end{axiom}

\begin{axiom}[A4 · U(1) / SU(2)]
Спинор $z\in\mathbb{C}^2$; gate collision --- SU(2) (\cref{ch:evolution}).
\end{axiom}

\begin{axiom}[A5 · Третий закон]
Абсолютный ноль $|z|=0$ недостижим; вакуум --- осциллирующий фон с $\rho_{\mathrm{field}}>0$.
\end{axiom}

\begin{axiom}[A6 · Второй закон (локально)]
Локальная фазовая энтропия на $\varepsilon$ не убывает при фиксированной норме $\Rightarrow$ mixing.
\end{axiom}

\begin{axiom}[A7 · Планковский потолок]
$|z|^2\le u_P$ (через $\rho_E/u_P\le 1$, знаменатель $+\varepsilon$).
\end{axiom}

\begin{axiom}[A8 · Макро-линейность]
При больших $|z|$ нелинейность затухает (гладкость на $\Tlayer$).
\end{axiom}

\begin{axiom}[A9 · Дискретная аналитичность]
$g$ удерживает дискретную голоморфность; CR-единственность $\Rightarrow$ глобальное поле на $\Lambda$.
\end{axiom}

\begin{axiom}[A10 · Топологический заряд]
Полюс на $\partial(\hv)$; $\oint d\arg=2\pi n$, $n\in\mathbb{Z}$.
\end{axiom}

\begin{axiom}[A11 · Устойчивость солитона]
Размазанное состояние не голоморфно; $K_P$ и $\Delta\phi$ --- затвор самофокусировки.
\end{axiom}

\begin{axiom}[A12 · Изотропия / Lorentz на $\Tlayer$]
Micro: isotropic streaming; macro-круг --- binomial readout + $\kgeom=1/\sqrt{2}$.
\end{axiom}

\begin{axiom}[A13 · Обратимость micro]
Leapfrog второго порядка на $\mathbb{Z}$; $\exists g^{-1}$ алгебраически.
\end{axiom}

\begin{axiom}[A14 · P / C / T / U(1)$_{\mathrm{vac}}$]
Дискретные симметрии на $g$ (\cref{ch:evolution}).
\end{axiom}

\begin{axiom}[A15 · Наименьшее действие]
$\kgeom$ из многогранника; $\gamma=\kappa_{\mathrm{link}}=1/|N|$; $\alpha^*=1+1/(4\pi)$.
\end{axiom}

\begin{axiom}[A16 · Паули]
Spin $1/2$: $2\pi\to -1$, $4\pi\to +1$; два совпадающих спинора в одном $\PlanckCell$ запрещены.
\end{axiom}

\begin{definition}[Глобальное состояние]
$z(x,t)\in\mathbb{C}^2$ на каждом $\PlanckCell$; $\Psi^t=\{z(x,t)\}_{x\in\Lambda}$;
$\sum_x z^\dagger z=\mathrm{const}$ (A3).
\end{definition}

\section{Детерминизм и «квантовая случайность''}

\begin{proposition}[Детерминизм $\Mlayer$]
\label{prop:determinism}
$g$ строго детерминирована: $z(x,t+1)=g(\{z(y,t)\}_{y\in N(x)})$ --- функция прошлого (A2).
RNG в фундаментальном шаге отсутствует.
\end{proposition}

Вероятность на $\Tlayer$ возникает из (i) динамического хаоса A5/A6 и (ii) грубого прибора:
binomial readout по $\gg 1$ ячейкам (\cref{ch:macro}).
$|\Psi|^2$ и Born --- статистика macro-фильтра, не кости в $g$.

\section{Теорема о невозможности тепловой смерти \texorpdfstring{$\Mlayer$}{M}}
\label{sec:theorem-23}

\begin{definition}[Shutdown · D1]
Существует $t$ без шага $g$, или $\htick\to\infty$, или $g$ меняется / обнуляется.
\end{definition}

\begin{definition}[M heat death · D2]
$\sum_x|z|^2\to 0$, или глобально $z\equiv 0$ устойчиво, или $g(\Psi^*)=\Psi^*$ с «мёртвым'' вакуумом без A5-кипения.
\end{definition}

\begin{definition}[$\Tlayer$ heat death · D3]
Macro-readout $\Phi=\mathcal{B}z$ теряет структуру ($\mathrm{Var}(\Phi)\to 0$) при $\sum|z|^2\approx\mathrm{const}$ на $\Mlayer$.
\end{definition}

\begin{remark}
$\mathcal{B}$ --- binomial readout $(1{:}2{:}1)^{\otimes 2R}$ (\cref{ch:macro}).
При $R\ge 1$ ядро $\mathcal{B}$ не инъективно.
\end{remark}

\begin{lemma}[Нет shutdown]
\label{lem:231}
При постулате 0.3 (\cref{ch:foundations}) для всех $t\in\mathbb{Z}$:
$\Psi^{t+1}=g(\Psi^t)$, $\htick$ и $g$ постоянны.
\end{lemma}

\begin{proof}
Постулат 0.3 определяет $\Mlayer$ как КА с одним $\htick$ на тик.
«Выключить'' --- нарушить постулат 0.3.
\end{proof}

\begin{lemma}[Сохранение информации]
\label{lem:232}
$N(t):=\sum_x z^\dagger z(x,t)$ --- инвариант: $N(t+1)=N(t)$ для всех $t$.
\end{lemma}

\begin{proof}
A3: $g=g_{\mathrm{nl}}\circ g_{\mathrm{lin}}$; $g_{\mathrm{lin}}$ --- product локальных unitaries на $N$;
$g_{\mathrm{nl}}$ --- $R(\Phi)=\exp(i\Phi)$ на $\mathbb{C}^2$, не меняет модуль.
Product unitaries сохраняют $N$ на каждом шаге. Индукция по $t$.
\end{proof}

\begin{corollary}
\label{cor:232a}
$N\to 0$ на $\Mlayer$ невозможно, если $N(0)>0$.
\end{corollary}

\begin{lemma}[A6 не отменяет A3]
\label{lem:233}
A6 не влечёт $N\to 0$.
\end{lemma}

\begin{proof}
A6 --- перераспределение фазы при фиксированной норме; A3 фиксирует $N$.
\end{proof}

\begin{lemma}[Нет устойчивого $z\equiv 0$]
\label{lem:234}
$z(x,t)\equiv 0$ не устойчивое термодинамическое состояние $g$.
\end{lemma}

\begin{proof}
A5: минимум $V(|z|^2)$ не в $z=0$; при $|z|\to 0$ gate $\Phi\to\mathrm{large}$.
При $z=0$ локальный deadlock ($z e^{i\Phi}=0$); A5 запрещает трактовать это как физический вакуум.
A5-согласованное начальное условие порождает permanent micro-кипение; $z\equiv 0$ недостижим
как глобальный attractor. Вместе с \cref{cor:232a}: D2 через $N\to 0$ или $z\equiv 0$ невозможен.
\end{proof}

\begin{lemma}[Micro-обратимость]
\label{lem:235}
$g$ на $Z_N[i]$ обратима: $\exists g^{-1}$, $g^{-1}(g(\Psi))=\Psi$.
\end{lemma}

\begin{proof}
A13: leapfrog $Z^++Z^-=2Z+\lfloor\mathcal{N}\rfloor$;
reverse $Z^-=2Z+\lfloor\mathcal{N}\rfloor-Z^+$. Kick-ledger + swap $(Z_t,Z_{t-1})$.
\end{proof}

\begin{lemma}[$\mathcal{B}$ необратима]
\label{lem:236}
При $R\ge 1$ существуют $z_1\ne z_2$ с $\mathcal{B}z_1=\mathcal{B}z_2$.
\end{lemma}

\begin{proof}
$\mathcal{B}$ --- свёртка с ядром $[1,2,1]/4$; у конечного $\Lambda$ ядро имеет нетривиальное ядро.
\end{proof}

\begin{proposition}[D3 не нарушает $\Mlayer$]
\label{prop:237}
$\nu_{\mathrm{CA}}$-затухание $\Phi$ на $\Tlayer$ --- не член $g$; $N$ на $\Mlayer$ invariant (\cref{lem:232}).
\end{proposition}

\begin{proof}
$\nu_{\mathrm{CA}}\nabla^2\Phi$ добавляется только в $\Tlayer$-уравнение после $\mathcal{B}$.
Macro drain --- перекачка в micro-кипение A5 при фиксированном $N$ (\cref{lem:233}).
\end{proof}

\begin{definition}[Frozen tick · D4]
$\lfloor\mathcal{N}(\Psi)\rfloor\equiv 0$ на всём $\Lambda$.
\end{definition}

\begin{definition}[Global fixed point · D5]
$g(\Psi^*)=\Psi^*$ при симметричной истории $(Z_t,Z_{t-1})=(Z^*,Z^*)$ $\Leftrightarrow$ D4 на $\Psi^*$.
\end{definition}

\begin{lemma}[Классификация D5]
\label{lem:238a}
D5 $\Leftrightarrow$ $z(x,t)\equiv c$ на всём торе $\Lambda$, $c\in Z_N[i]$ (включая $c=0$).
\end{lemma}

\begin{proof}
($\Rightarrow$) Symmetric leapfrog при $Z^+=Z^-=Z^*$ $\Rightarrow$ $\lfloor\mathcal{N}\rfloor=0$
$\Rightarrow$ $\Phi=0$ $\Rightarrow$ $\zeta_{\mathrm{imag}}=0$. Discrete CR на каждой ячейке;
$\Delta_{\mathrm{disc}} z=0$ $\Leftrightarrow$ $z(x)=\langle z\rangle_N$.
Maximum principle на связном графе $\Rightarrow$ $z(x)\equiv c$.

($\Leftarrow$) $z\equiv c$ $\Rightarrow$ $\Phi=0$, $\lfloor\mathcal{N}\rfloor=0$, тождество leapfrog.
\end{proof}

\begin{lemma}[A5/A6 исключают физические D5]
\label{lem:238b}
Ни одна конфигурация из \cref{lem:238a} не входит в A5-физическое состояние.
\end{lemma}

\begin{proof}
$c=0$: deadlock; A5 --- $z\equiv 0$ не вакуум.
$c\ne 0$ static: A5 требует кипящий фон; A6 запрещает «замёрзшие'' фазы.
\end{proof}

\begin{theorem}[Нет A5-физических fixed points без кипения]
\label{thm:238}
При §0.3, A3, A5, A6, A9, A13 и $g$ §3.12 полный класс A5-физических global fixed points
$g(\Psi^*)=\Psi^*$ без micro-кипения --- пуст.
\end{theorem}

\begin{proof}
\cref{lem:238a,lem:238b}.
\end{proof}

\begin{theorem}[Главная теорема §2.3]
\label{thm:main-23}
При §0.3, A3, A5, A6, A13 и $g$ §3.12:
\begin{enumerate}[label=(\roman*)]
\item $\neg$D1 (нет shutdown) --- \cref{lem:231};
\item $\neg$D2 (нет M heat death) --- \cref{cor:232a,lem:234,thm:238};
\item A6 и A3 согласованы --- \cref{lem:233};
\item micro обратим, macro необратим --- \cref{lem:235,lem:236,prop:237};
\item нет A5-физических D5 --- \cref{thm:238}.
\end{enumerate}
\end{theorem}

\begin{proof}
Прямо из перечисленных лемм и теорем.
\end{proof}

\begin{remark}
Модельная вселенная не «выключается'' и не «умирает'' в смысле D2;
«тепловая смерть'' наблюдателя --- D3 на $\Tlayer$.
Космологическое начало наблюдаемой вселенной --- meta-слой, не первый тик $\Mlayer$.
\end{remark}
